% !TEX program = xelatex
% 保险行业季度调研报告生成器 — Beamer 宣传演示
% 编译：xelatex 保险报告生成器_Beamer宣传.tex （编译两遍）
\documentclass[aspectratio=169,10pt,UTF8]{ctexbeamer}

\usetheme{Madrid}
\usecolortheme{whale}
\usefonttheme{professionalfonts}
\setbeamertemplate{navigation symbols}{}
\setbeamertemplate{footline}[frame number]

% 品牌色（与桌面软件深蓝风格一致）
\definecolor{BrandBlue}{HTML}{17365D}
\definecolor{BrandLight}{HTML}{EAF0F7}
\definecolor{BrandAccent}{HTML}{2E75B6}
\definecolor{BrandGreen}{HTML}{375623}
\definecolor{BrandWarn}{HTML}{C00000}

\setbeamercolor{structure}{fg=BrandBlue}
\setbeamercolor{palette primary}{bg=BrandBlue,fg=white}
\setbeamercolor{palette secondary}{bg=BrandAccent,fg=white}
\setbeamercolor{palette tertiary}{bg=BrandBlue,fg=white}
\setbeamercolor{title}{fg=BrandBlue}
\setbeamercolor{frametitle}{bg=BrandBlue,fg=white}
\setbeamercolor{block title}{bg=BrandBlue,fg=white}
\setbeamercolor{block body}{bg=BrandLight,fg=black}
\setbeamercolor{block title alerted}{bg=BrandWarn,fg=white}
\setbeamercolor{block title example}{bg=BrandGreen,fg=white}
\setbeamercolor{item}{fg=BrandAccent}

\usepackage{booktabs}
\usepackage{tabularx}
\usepackage{graphicx}
\usepackage{tikz}
\usetikzlibrary{positioning}
\usepackage{hyperref}
\hypersetup{colorlinks=true,linkcolor=BrandBlue,urlcolor=BrandAccent}

\title{\textbf{保险行业季度调研报告生成器}}
\subtitle{多源 DOCX 一键合成 · 智能章节 · 可审可调 · 终版交付}
\author{保险行业研究工具}
\institute{桌面端 · Windows / macOS}
\date{2026}

\begin{document}

% =====================================================================
\begin{frame}[plain]
  \vfill
  \begin{center}
    {\color{BrandBlue}\Huge\bfseries 保险行业季度调研报告生成器}\\[0.6em]
    {\Large 把多部门材料，变成一份可交付的季度终版}\\[1.2em]
    \begin{tikzpicture}
      \node[fill=BrandLight,rounded corners=6pt,inner sep=10pt,text width=0.82\textwidth,align=center] {
        \large 输入多份 DOCX \quad$\rightarrow$\quad 预览与调序 \quad$\rightarrow$\quad 导出 DOCX / PDF
      };
    \end{tikzpicture}\\[1.4em]
    {\normalsize 桌面图形界面 · 格式模板复用 · 跨平台}
  \end{center}
  \vfill
  \begin{flushright}
    \small 2026 \quad|\quad 保险行业研究工具
  \end{flushright}
\end{frame}

% =====================================================================
\begin{frame}{为什么需要它？}
  \begin{columns}[T]
    \begin{column}{0.48\textwidth}
      \begin{block}{传统方式的痛点}
        \begin{itemize}
          \item 各部门各自写一节，手工粘贴合并
          \item 标题层级、编号、字体不统一
          \item 目录、题注、页码反复返工
          \item 季度切换时结构难对齐历史终版
          \item 一处改编号，全文引用易出错
        \end{itemize}
      \end{block}
    \end{column}
    \begin{column}{0.48\textwidth}
      \begin{exampleblock}{用生成器之后}
        \begin{itemize}
          \item \textbf{多文件自动汇总} 到统一报告
          \item \textbf{章节识别 + 预览编辑} 可控
          \item \textbf{格式模板} 对齐上季终版
          \item 一键导出 \textbf{DOCX / PDF}
          \item 守恒校验，减少静默丢内容
        \end{itemize}
      \end{exampleblock}
    \end{column}
  \end{columns}
  \vspace{0.8em}
  \centering
  \textit{「材料还是各部门写，终版由工具按规则组装。」}
\end{frame}

% =====================================================================
\begin{frame}{一句话定位}
  \vspace{1.2em}
  \begin{center}
    \begin{tikzpicture}
      \node[
        fill=BrandBlue, text=white, rounded corners=10pt,
        inner sep=18pt, text width=0.88\textwidth, align=center
      ] {
        \Large\bfseries
        面向保险行业季度调研场景的\\[0.35em]
        多源 Word 报告自动合成桌面工具
      };
    \end{tikzpicture}
  \end{center}
  \vspace{1.2em}
  \begin{columns}[T]
    \begin{column}{0.32\textwidth}
      \centering
      {\color{BrandAccent}\Large\bfseries 合}\\[0.3em]
      多份输入 DOCX\\合并为结构清晰的终稿
    \end{column}
    \begin{column}{0.32\textwidth}
      \centering
      {\color{BrandAccent}\Large\bfseries 控}\\[0.3em]
      章节预览、改名、\\调序、改层级
    \end{column}
    \begin{column}{0.32\textwidth}
      \centering
      {\color{BrandAccent}\Large\bfseries 齐}\\[0.3em]
      字体、页眉、目录\\与上季终版对齐
    \end{column}
  \end{columns}
\end{frame}

% =====================================================================
\begin{frame}{核心能力一览}
  \begin{itemize}
    \item \textbf{智能章节识别}：支持「第N章 / 2.1 / 2.1.1」等多级编号与 Word 标题样式
    \item \textbf{同号合并}：多文件同 section 自动汇总，冲突可在预览中发现
    \item \textbf{人工可控}：预览表中改标题、编号、层级，并支持拖拽式调序
    \item \textbf{格式模板}：可选历史上季终版，复用章节顺序、标题与字体风格
    \item \textbf{图文保真}：表格、图片、加粗等随内容迁移
    \item \textbf{题注与引用}：图/表题注与交叉引用尽量在合成时理顺
    \item \textbf{一键交付}：输出 DOCX；本机有 Word / LibreOffice 时可出 PDF
  \end{itemize}
\end{frame}

% =====================================================================
\begin{frame}{三步出报告}
  \begin{enumerate}
    \item[\textcolor{BrandBlue}{\textbf{1}}]
      \textbf{添加材料}\\[0.2em]
      {\small 导入本季度各部门 DOCX；可选加载「上季终版」作格式模板}
      \vspace{0.55em}
    \item[\textcolor{BrandBlue}{\textbf{2}}]
      \textbf{预览与修订}\\[0.2em]
      {\small 查看识别出的章节树；修改标题/编号/层级；调整顺序}
      \vspace{0.55em}
    \item[\textcolor{BrandBlue}{\textbf{3}}]
      \textbf{生成交付}\\[0.2em]
      {\small 填写报告周期、标题、单位、页眉；导出终版 DOCX（及 PDF）}
  \end{enumerate}
  \vspace{0.8em}
  \begin{block}{典型场景}
    用 \textbf{2026Q2 终版} 作模板，放入 \textbf{2026Q3 六份输入}，预览确认后一键生成第三季度终版。
  \end{block}
\end{frame}

% =====================================================================
\begin{frame}{界面与参数（桌面端）}
  \begin{columns}[T]
    \begin{column}{0.52\textwidth}
      \begin{itemize}
        \item 原生图形界面（Tk）
        \item 左侧：模板、报告参数、页眉、字体
        \item 右侧：文件列表、章节预览表、日志
        \item 预览后才允许「生成」，避免误点
        \item 生成完成后可直接打开结果
      \end{itemize}
    \end{column}
    \begin{column}{0.45\textwidth}
      \begin{exampleblock}{可配置项}
        \small
        \begin{itemize}
          \item 周期 / 主副标题 / 单位 / 日期
          \item 页眉：标题·单位·日期·周期
          \item 一/二/三级标题与正文字体字号
          \item 正文首行缩进
          \item 是否启用模板标题
        \end{itemize}
      \end{exampleblock}
    \end{column}
  \end{columns}
\end{frame}

% =====================================================================
\begin{frame}{质量与可靠性}
  \begin{columns}[T]
    \begin{column}{0.48\textwidth}
      \begin{block}{内容守恒}
        \begin{itemize}
          \item 输入块数与写出块数校验
          \item 坏文件隔离，不拖垮整批
          \item 合并冲突在预览中可见
        \end{itemize}
      \end{block}
      \vspace{0.5em}
      \begin{block}{包完整性}
        \begin{itemize}
          \item 生成后校验 DOCX 内部关系
          \item 图片/超链接尽量保真复制
        \end{itemize}
      \end{block}
    \end{column}
    \begin{column}{0.48\textwidth}
      \begin{alertblock}{智能防误识别}
        正文中的「\textbf{10年期国债}」「\textbf{7月保费}」「\textbf{2026年…}」
        \textbf{不会}再被误判成
        「第十章 / 第七章 / 第2026章」。
        \vspace{0.4em}

        {\small 一级编号须有分隔符；过滤年份号与「N年/N月」粘连正文。}
      \end{alertblock}
    \end{column}
  \end{columns}
\end{frame}

% =====================================================================
\begin{frame}{格式模板：季度间「长得像」}
  \centering
  \begin{tikzpicture}[
    box/.style={draw=BrandBlue, rounded corners, fill=BrandLight,
      align=center, inner sep=8pt, text width=2.6cm, minimum height=1.5cm},
    arr/.style={->, thick, BrandAccent}
  ]
    \node[box] (q2) {\textbf{Q2 终版}\\格式模板};
    \node[box, right=1.6cm of q2] (in) {\textbf{Q3 输入}\\多份 DOCX};
    \node[box, right=1.6cm of in] (out) {\textbf{Q3 终版}\\结构对齐};
    \draw[arr] (q2) -- (in);
    \draw[arr] (in) -- (out);
  \end{tikzpicture}

  \vspace{1em}
  \begin{itemize}
    \item 模板提供：\textbf{章节顺序}、\textbf{标题措辞}、\textbf{字体层级风格}
    \item 输入提供：本季最新正文、图表与数据
    \item 可开关「匹配章节时优先使用模板标题」
    \item 仍保留人工预览修订权，模板不是黑盒
  \end{itemize}
\end{frame}

% =====================================================================
\begin{frame}{跨平台交付}
  \begin{table}
    \centering
    \renewcommand{\arraystretch}{1.25}
    \begin{tabular}{>{\bfseries}l ll}
      \toprule
      平台 & 使用方式 & PDF \\
      \midrule
      Windows & 双击 \texttt{保险报告生成器.exe} & 本机 Word（推荐） \\
      macOS & \texttt{run\_mac.command} 或打包 .app & Word / LibreOffice \\
      源码 & \texttt{工具/} + Python 依赖 & 同上多后端 \\
      \bottomrule
    \end{tabular}
  \end{table}
  \vspace{0.6em}
  \begin{block}{说明}
    \begin{itemize}
      \item Windows 安装包与 \texttt{\_internal} 需同目录放置
      \item Mac 无法直接运行 \texttt{.exe}；请用源码或在 Mac 上打 \texttt{.app}
      \item 无 Word 时仍可稳定导出 \textbf{DOCX}，PDF 可后补
    \end{itemize}
  \end{block}
\end{frame}

% =====================================================================
\begin{frame}{适合谁用？}
  \begin{itemize}
    \item \textbf{研究团队}：每季度多部门供稿、需统一终版
    \item \textbf{报告统筹人}：要控结构、标题与交付节奏
    \item \textbf{分析师}：专注写自己的章节，不必手工排全书
    \item \textbf{质控同事}：预览阶段即可发现冲突与异常章节
  \end{itemize}
  \vspace{0.8em}
  \begin{exampleblock}{不替代什么}
    不替代业务判断与文字写作；替代的是\textbf{合并、编号、格式对齐、反复粘贴}这类机械劳动。
  \end{exampleblock}
\end{frame}

% =====================================================================
\begin{frame}{价值总结}
  \begin{columns}[T]
    \begin{column}{0.33\textwidth}
      \centering
      {\color{BrandBlue}\Large\bfseries 更快}\\[0.4em]
      从「手工拼一周」\\到「预览后一键出」
    \end{column}
    \begin{column}{0.33\textwidth}
      \centering
      {\color{BrandBlue}\Large\bfseries 更稳}\\[0.4em]
      防伪章节、守恒校验\\坏文件隔离
    \end{column}
    \begin{column}{0.33\textwidth}
      \centering
      {\color{BrandBlue}\Large\bfseries 更齐}\\[0.4em]
      模板对齐历史终版\\字体页眉统一
    \end{column}
  \end{columns}
  \vspace{1.4em}
  \begin{center}
    \begin{tikzpicture}
      \node[fill=BrandLight,rounded corners,inner sep=12pt,text width=0.9\textwidth,align=center] {
        \large\textbf{同一套流程，服务每一个保险行业调研季度。}
      };
    \end{tikzpicture}
  \end{center}
\end{frame}

% =====================================================================
\begin{frame}{立即开始}
  \begin{enumerate}
    \item 打开文件夹中的 \textbf{保险报告生成器.exe}（Windows）
    \item 或进入 \textbf{工具/}，按 \texttt{Mac使用说明.md} 在 Mac 运行
    \item 准备本季输入 DOCX；可选指定上季终版为模板
    \item 预览 $\rightarrow$ 微调 $\rightarrow$ 生成
  \end{enumerate}
  \vspace{0.9em}
  \begin{block}{资源位置（示例）}
    \small
    \begin{itemize}
      \item 程序与样例：本地「报告合成」目录
      \item 源码：\texttt{工具/desktop\_app.py}、\texttt{report\_aggregator.py}
      \item 远程：\url{https://github.com/JimLi0128/AI-MunichRe/tree/mzj}
    \end{itemize}
  \end{block}
\end{frame}

% =====================================================================
\begin{frame}[plain]
  \vfill
  \begin{center}
    {\color{BrandBlue}\Huge\bfseries 谢谢}\\[0.8em]
    {\Large 保险行业季度调研报告生成器}\\[0.5em]
    {\large 多源合成 · 章节可控 · 终版可交付}\\[1.4em]
    \begin{tikzpicture}
      \node[fill=BrandBlue,text=white,rounded corners=8pt,inner sep=12pt] {
        \large 问题与反馈欢迎随时提出
      };
    \end{tikzpicture}
  \end{center}
  \vfill
\end{frame}

\end{document}
